\section{Related Work}
Intermittent systems research addresses the consistency of execution across
power failures. Alpaca organizes a program into tasks, privatizes shared data
and commits task effects so interrupted execution can restart without exposing
partial updates~\cite{alpaca}. This establishes a concrete recovery mechanism;
it does not make a current energy summary sufficient for a different capacitor
or establish that an event reached its receiver. Our retained-state witnesses
use an idealized latch and do not implement or evaluate Alpaca.

Curricle specifies intermittent idempotence requirements through a type system
and information-flow qualifiers, checking programs against those
requirements~\cite{curricle}. Its program-level consistency question complements
our admission question: whether an observation has declared units and domain,
is independently qualified on the supplied source, and reaches a compatible
design consumer. Passing these interface checks does not prove intermittent
idempotence, arbitrary program correctness or physical validity.

Counterexample-guided abstraction refinement (CEGAR) refines an abstraction
when an abstract counterexample is spurious in the concrete
system~\cite{cegar}. Our formulation shares the use of distinguishability to
explain missing information. Its role is narrower: two legal digital worlds
with the same representation require different verdicts after a specified
intervention. Proposition~\ref{prop:factorization} is an elementary
factorization condition, not a new general theory of abstraction or a CEGAR
algorithm. We supply neither a general model checker nor an automated
refinement guarantee. The empirical contribution is the source-bound
comparison among complete static information, executable feedback and
qualified live observations, including failed chains and receiver-side
service conditions.
